\documentclass[12pt]{article}
\usepackage[T2A]{fontenc}
\usepackage[utf8]{inputenc}
\usepackage[russian]{babel}
\usepackage{amsmath,amssymb}
\usepackage{geometry}
\geometry{margin=2.5cm}
\usepackage{indentfirst}

\title{Магнитные скирмионы в ван-дер-ваальсовых гетероструктурах:\\ управление топологией электрическим полем}
\author{Х.~ван дер Берг}
\date{2024}

\begin{document}
\maketitle

\section*{Аннотация}
Показано, что в гетероструктурах на основе двумерных магнитных материалов с нарушенной инверсионной симметрией топологическая фаза скирмионной решётки может переключаться внешним электрическим полем. Построена эффективная спиновая модель с учётом взаимодействия Дзялошинского–Мории и магнитоэлектрического эффекта.

\vspace{0.5cm}
\noindent\textbf{Ключевые слова:} скирмионы; ван-дер-ваальсовы гетероструктуры; магнитоэлектрический эффект; взаимодействие Дзялошинского–Мории; энергонезависимая память.

\vspace{0.3cm}
\noindent\textbf{Раздел журнала:} Физика. \quad \textbf{Файл:} 1 из 2.

\section{Введение}
Скирмионы — топологически защищённые магнитные структуры, перспективные для устройств памяти. Управление ими электрическим полем остаётся открытой задачей.

\section{Модель}
Эффективный гамильтониан:
\begin{equation}
\mathcal{H} = -J \sum_{\langle ij \rangle} \mathbf{S}_i \cdot \mathbf{S}_j + \sum_{\langle ij \rangle} \mathbf{D}_{ij} \cdot (\mathbf{S}_i \times \mathbf{S}_j) - \mathbf{E} \cdot \mathbf{P}.
\end{equation}

\section{Заключение}
Предложена схема энергонезависимой памяти на одиночных скирмионах.

\end{document}